\section{Methodology}
\label{sec:method}

We describe our approach for testing whether MLP layers can be approximated by shallow programs over token directions.
The core idea is to project MLP inputs and outputs into vocabulary space via the unembedding matrix and measure how well a linear model in this space approximates the MLP's behavior.

\subsection{Setup}
\label{sec:setup}

\para{Model.}
We use \gpttwo (12 layers, $\dmodel = 768$, $\dmlp = 3072$, vocabulary size $|\gV| = 50{,}257$), loaded via TransformerLens~\citep{nanda2022transformerlens}.
GPT-2 Small is the most widely studied model in the mechanistic interpretability literature, enabling direct comparison with prior work.

\para{Data.}
We use the \wikitext validation split~\citep{merity2017pointer}, consistent with the datasets used by \citet{geva2021transformer,geva2022transformer}.
We sample 200 text sequences, each truncated to 128 tokens (with BOS token prepended), filtering sequences shorter than 10 tokens.
This yields 19,281 token positions across all sequences.
At each layer $\ell$ and token position $t$, we collect the MLP input $\vx_{\ell,t} \in \sR^{\dmodel}$ (post-LayerNorm) and MLP output $\vy_{\ell,t} \in \sR^{\dmodel}$.

\subsection{Token Direction Decomposition}
\label{sec:decomposition}

We project MLP activations into vocabulary space using the unembedding matrix $\Wu \in \sR^{|\gV| \times \dmodel}$.
For an activation vector $\vx \in \sR^{\dmodel}$, the token logits are:
\begin{equation}
    \vz = \Wu \vx \in \sR^{|\gV|}
    \label{eq:projection}
\end{equation}
where the $i$-th component $z_i$ measures how strongly $\vx$ aligns with the direction of token $i$ in the unembedding space.

To measure how well token directions capture MLP activations, we select the top-$k$ tokens (by logit magnitude) for each vector, extract the corresponding rows of $\Wu$ as a basis $\mB_k \in \sR^{k \times \dmodel}$, reconstruct via least-squares projection $\hat{\vx} = \mB_k^{\dagger} \mB_k \vx$, and measure cosine similarity and $\Rsq$ between $\vx$ and $\hat{\vx}$.

\subsection{Shallow Program Approximation}
\label{sec:shallow}

The central experiment fits a linear model mapping input token logits to output token logits:
\begin{equation}
    \hat{\vz}^{\text{out}} = \mA \, \vz^{\text{in}} + \vb
    \label{eq:shallow}
\end{equation}
where $\vz^{\text{in}} = \Wu \vx \in \sR^{|\gV|}$ and $\vz^{\text{out}} = \Wu \vy \in \sR^{|\gV|}$ are the input and output token logits, $\mA \in \sR^{k \times k}$ is a learned weight matrix, and $\vb$ is a bias term.

To keep the model tractable, we select the top-$k$ most variable token dimensions (across all positions) for both inputs and outputs.
We vary $k \in \{10, 25, 50, 100, 200, 500\}$ to study the effect of the number of token directions.
We fit Ridge regression ($\alpha = 1.0$) on a 70/30 train/test split and report $\Rsq$, cosine similarity, and top-5 agreement on the held-out test set.

\subsection{Baselines}
\label{sec:baselines}

We compare the token-direction program against two baselines:

\para{Random directions.}
We replace $\Wu$ with a random orthogonal matrix $\mR \in \sR^{k \times \dmodel}$ (rows drawn from a random orthogonal basis) and fit the same linear model.
If the token-direction program succeeds because the vocabulary is a privileged basis, it should outperform random directions.
If random directions match or exceed token directions, the success is attributable to the MLP's inherent linearity.

\para{Full linear map.}
We fit a linear map directly in the residual stream: $\hat{\vy} = \mA' \vx + \vb'$ where $\mA' \in \sR^{\dmodel \times \dmodel}$.
This provides an upper bound on the variance explainable by any linear model, without the information bottleneck of projecting through token space.

\subsection{Additional Analyses}
\label{sec:additional}

\para{MLP linearity.}
We measure how well a full-rank $\dmodel \to \dmodel$ linear map approximates each MLP layer, providing a direct characterization of per-layer nonlinearity.

\para{Activation sparsity.}
We analyze the sparsity of MLP hidden activations (post-GELU) to test whether computation is concentrated in a few neurons.
We report the fraction of neurons with $|\text{activation}| > 0.1$ and the cumulative contribution of the top-$n$ neurons.

\para{Next-token prediction.}
We project MLP outputs to vocabulary space and measure how often the top-ranked token matches the actual next token, assessing whether individual MLP layers directly predict the next token.
